problem:  
Assume there exists a well‑defined notion of a **4‑dimensional singular Ricci flow**, obtained as a pointed Cheeger–Gromov limit of admissible Ricci flows with surgery whose surgery parameters \( \delta_k \to 0 \), in direct analogy with the Bamler–Kleiner construction in dimension three.  

Let \( (M^4, g_0) \) be a smooth, closed, oriented Riemannian \(4\)-manifold.  For each admissible surgery scheme, let \( (M^4, g_k(t)) \) denote the resulting Ricci flow with surgery starting at \(g_0\) and existing for all \(t \ge 0\).  
Suppose that along each sequence \( (g_k(t)) \) the space–time metrics subconverge (as \( \delta_k \to 0 \)) to a limiting singular flow \( (M^4, g_\infty(t)) \).  

**Question.** — *Assuming this framework exists,* is the resulting limit \( (M^4, g_\infty(t)) \) **unique**—that is, independent (up to diffeomorphism and time reparametrization) of the particular surgery scheme and neck‑scale parameters?  
If not, can one construct examples (or identify obstructions) where distinct families of such surgery flows for the same initial \( (M^4, g_0) \) yield analytically inequivalent limit singular flows?  

Why it is a "good" problem:  
This conditional problem acknowledges the current state of the art—namely, that a rigorous 4‑dimensional singular Ricci flow theory analogous to the established 3‑dimensional case has not yet been fully built. By first hypothesizing its existence, the question isolates the deep issue of **analytic uniqueness and canonical evolution** in higher‑dimensional Ricci flow.  
It probes the boundary between geometry and analysis: if uniqueness fails, the smooth topology of four‑manifolds might support multiple inequivalent “weak continuations” of Ricci flow past singularities; if uniqueness holds, it would suggest a powerful extension of Perelman‑type canonical geometry into higher dimensions.  
The statement is compact and conceptually natural—"If a 4‑D singular Ricci flow exists, is it unique?"—but the implications reach profound unsolved territories in geometric analysis, topology, and the theory of weak flows. This makes it a truly research‑level “narrow entrance, profound depth” problem.